\section{Dated observations and a common decision rule}
\label{sim:revised-source-actor-rules}

The analyst's receipt assumption does not tell us what the manager can observe. Two explicit alternatives make the distinction testable. The first retains the original numeric observations: the first coupon signal and the later trigger and cash-coverage signal. The second pairs each numeric message with the receipt reported by the financial assertion selected at that date. These are supplied information assumptions. A reported future receipt is neither cash already received nor proof of the eventual amount.

For a correction arriving between decisions, the supplemented first message reports 7 and the second reports 6, even when the analyst evaluates a complete receipt-6 world. Putting the later correction into the first message would give the manager information too early. At the second decision, the key still includes the first message and the manager's first action. That memory is needed even if the current numeric signal coincides with one reached after a different earlier choice.

Several supplied source/world alternatives can then be evaluated together. The alternatives share an actor, dates and action meanings. Their identities remain internal: equal visible decision keys require the same action in every alternative. Finding a policy separately in each world is insufficient. For example, if one world requires pay and another requires draw at the same first message, each has a policy in isolation but there is no common rule. The retained economic-world declaration still requires a correction ancestry compatible with its supplied world; this comparison pools independently specified alternatives rather than inferring several truths from one corrected report.

The prospective domain is formed before selecting favorable source outcomes, using the original source heads and the declared complete worlds. Adding the reported receipt can expand this domain beyond the twelve-key enumeration allowance. The larger domain is retained intact. A size-limit result says that the complete policy population has not been enumerated; it says nothing about whether a policy exists. The original same-world observation-correction example retains exactly its 16 policies, and pairing its messages with a constant reported receipt preserves that number.\footnote{\path{data/accounting-master/revised-source-actor-messages-20261027/attempt-02/} preserves 47 comparisons, 394 assertions and 20 alteration controls. Independent reconstruction checks the dated reports, complete first-stage memory, protected policy population and shared-action restrictions. Unresolved sources and worlds retain their separate statuses.}

\subsection{An action assignment or a contradiction can settle existence}

The two available actions give a smaller logical question than enumerating every rule. Associate a Boolean variable $x_i$ with each complete decision key: 0 denotes pay and 1 denotes draw. A financial row refused at the first decision forbids one assignment $x_i=a$. A row refused at the second decision forbids the joint choice of its first and second actions, $x_i=a$ and $x_j=b$. A rule avoids those refusals exactly when it satisfies, respectively,
\begin{equation}
 x_i\ne a,
 \qquad
 (x_i\ne a)\ \lor\ (x_j\ne b).
 \label{eq:policy-refusal-clauses}
\end{equation}
The second clause says that at least one of the actions leading to the refused row must change. It is equivalent to the implications
\begin{equation}
 (x_i=a)\Rightarrow(x_j=1-b),
 \qquad
 (x_j=b)\Rightarrow(x_i=1-a).
 \label{eq:policy-refusal-implications}
\end{equation}
A unary refusal similarly implies $(x_i=a)\Rightarrow(x_i=1-a)$. These clauses preserve the original financial rows and observation keys. They change how existence is decided, without changing which decisions are allowed.

A proposed total assignment is checked directly against every refused row. If none is reached under that assignment, it supplies one common policy for the resolved financial family. It does not enumerate other policies or choose an optimal one. Conversely, suppose implications justified by refused rows lead from $x_i=0$ to $x_i=1$ and also from $x_i=1$ to $x_i=0$. Either value contradicts a consequence of choosing it. The two paths therefore certify that no assignment can avoid all those refusals. A graph calculation can discover the paths, but verification checks each edge against its cited financial row.

In the between-date receipt-6 example, a three-edge contradiction has a direct decision interpretation. At the first numeric message X paired with reported receipt 7, paying is refused, so the rule must draw. After that action, a later message reports no trigger, inadequate coverage and receipt 6. Paying is refused there, and drawing is also refused. The two second-stage refusals make the required first draw impossible. This identifies incompatible requirements in the supplied decision problem. It does not yet tell us whether a particular economic path is impossible even with full information, or whether indistinguishable paths demand different actions.

The certificate comparison decides thirteen domains above twelve keys, including a satisfying rule for a 30-key reverse-receipt declaration and contradictions for receipt-6 declarations with 38 keys and a concurrent-correction declaration with 54 keys. The previous enumeration results remain unchanged. Missing world rows receive no certificate. A contradiction found within a resolved subset refutes that subset, while a satisfying assignment there cannot establish feasibility for unresolved alternatives.\footnote{\path{data/accounting-master/policy-existence-certificates-20261027/} checks all 47 retained source/world comparisons and 260 small Boolean formulas against exhaustive assignments, with 1,083 assertions and 22 alteration controls. The 31 certificates comprise six assignments and 25 contradictions; two contradictions apply only to resolved subsets. Each small-domain result agrees with the protected full policy population.}

These results make the next question more specific. A rejected common rule calls for examining the financial histories responsible for its contradiction. That diagnosis must distinguish economic infeasibility from an observation that hides a consequential difference before proposing a repair to the model or its information assumptions.
